% CODEX: wstaw bezpośrednio przed \section{Pochodna kowariantna i transport wzdłuż krzywej}
% w 12-koneksje.tex. Istniejącego wzoru Koszula nie przepisuj.
\subsection*{Wzór Koszula w ramie nieholonomicznej: geometria Nil}
\label{sec:thur-dop-koszul-rama}
We współrzędnej ramie $\partial_i$ zachodzi $[\partial_i,\partial_j]=0$.
W geometrycznych modelach grup Liego znacznie wygodniejsze bywają jednak
ramy ortonormalne, których pola nie komutują. Nie wolno wtedy stosować
wzorów na symbole Christoffela tak, jakby były ramą współrzędną.

\begin{stwierdzenie}[Koneksja w ramie ortonormalnej]
\label{prop:thur-dop-koszul-rama}
Niech $E_1,\ldots,E_n$ będzie lokalną ramą ortonormalną,
a $c_{ij}^{k}=g([E_i,E_j],E_k)$. Wtedy dla koneksji Leviego-Civity
\begin{equation}\label{eq:thur-dop-koszul-rama}
2g(\nabla_{E_i}E_j,E_k)
=c_{ij}^{k}-c_{jk}^{i}+c_{ki}^{j}.
\end{equation}
\end{stwierdzenie}
\begin{proof}
Wstawiamy $E_i,E_j,E_k$ do wzoru Koszula
\eqref{eq:koszul}. Wszystkie pochodne $E_i(g(E_j,E_k))$
znikają, ponieważ $g(E_j,E_k)=\delta_{jk}$ jest stałe.
Pozostają trzy wyrazy zawierające nawiasy pól,
co daje powyższą formułę. Równość ta nie wymaga,
aby rama pochodziła od mapy współrzędnych.
\end{proof}

\begin{przyklad}[Koneksja Leviego-Civity w grupie Heisenberga]
\label{prz:thur-dop-koneksja-nil}
Dla metryki $g_N$ z Przykładu~\ref{prz:thur-dop-metryki-nil-sol}
weźmy ramę
$E_1=\partial_x$, $E_2=\partial_y+x\partial_z$,
$E_3=\partial_z$. Jest ona ortonormalna, a jedynym
niezerowym nawiasem bazowym (z dokładnością do antysymetrii)
jest $[E_1,E_2]=E_3$. Stąd~\eqref{eq:thur-dop-koszul-rama}
daje
\[
\begin{aligned}
 \nabla_{E_1}E_2&=\tfrac12E_3,&
 \nabla_{E_2}E_1&=-\tfrac12E_3,\\
 \nabla_{E_1}E_3&=-\tfrac12E_2,&
 \nabla_{E_3}E_1&=-\tfrac12E_2,\\
 \nabla_{E_2}E_3&=\tfrac12E_1,&
 \nabla_{E_3}E_2&=\tfrac12E_1.
\end{aligned}
\]
Ponadto $\nabla_{E_i}E_i=0$ dla każdego $i$.
Przykładowo $2g(\nabla_{E_1}E_2,E_3)
=g([E_1,E_2],E_3)=1$.
Różnica $\nabla_{E_1}E_2-\nabla_{E_2}E_1=E_3$
równa się $[E_1,E_2]$; koneksja nie ma więc torsji,
mimo że współczynniki w tej ramie nie są symetryczne.
Ten sam rachunek posłuży do wyznaczenia tensora krzywizny
w Przykładzie~\ref{prz:thur-dop-krzywizna-nil}.
\end{przyklad}
